\subsubsection{缓增分布的 Fourier 变换}\label{subsec:tempered-fourier-transform}

前一个定理给出了定义缓增分布的 Fourier 变换的动机.

\begin{Definition}
\label{def:tempered-fourier-transform}
设 $f\in\mathcal{S}^{\prime}(\mathbb{R}^m)$. 那么 $f$ 的\textbf{Fourier 变换}通过泛函定义为
\begin{equation}\label{eq:tempered-distribution-fourier-definition}
(\mathcal{F}[f],\phi)=(f,\mathcal{F}^{-1}[\phi]),\quad\phi\in\mathcal{S}(\mathbb{R}^m).
\end{equation}

\end{Definition}

由定义可知, $\mathcal{F}$ 是从 $\mathcal{S}^{\prime}(\mathbb{R}^m)$ 到其自身的连续映射. 也容易验证, 公式 \eqref{eq:fourier-frequency-derivatives} 和 \eqref{eq:fourier-spatial-derivatives} 在 $\mathcal{S}^{\prime}(\mathbb{R}^m)$ 中仍成立; 逆变换公式 \eqref{eq:fourier-inversion-formula} 也如此, 它可重新表述为
\begin{equation}\label{eq:fourier-transform-twice-reflection}
\mathcal{F}\mathcal{F}[f(\bm{x})]=f(-\bm{x});
\end{equation}
这一形式对广义函数有意义.

现在考察若干例子.

\begin{Example}
\label{ex:fourier-delta}
有
\begin{equation}\label{eq:fourier-transform-delta-pairing}
\begin{aligned}
(\mathcal{F}[\delta],\phi)&=(\delta,\mathcal{F}^{-1}[\phi])=\mathcal{F}^{-1}[\phi](0)\\
&=(2\pi)^{-m/2}\int_{\mathbb{R}^m}\phi(\bm{x})\,\mathrm{d}\bm{x},
\end{aligned}
\end{equation}
即 $\delta$ 的 Fourier 变换为常数 $(2\pi)^{-m/2}$.
\end{Example}

\begin{Example}
\label{ex:fourier-constant}
有
\begin{equation}\label{eq:fourier-transform-constant-pairing}
\begin{aligned}
(\mathcal{F}[1],\phi)&=(1,\mathcal{F}^{-1}[\phi])=\int_{\mathbb{R}^m}\mathcal{F}^{-1}[\phi](\bm{x})\,\mathrm{d}\bm{x}\\
&=(2\pi)^{m/2}\mathcal{F}\mathcal{F}^{-1}[\phi](0)=(2\pi)^{m/2}\phi(0),
\end{aligned}
\end{equation}
即 1 的 Fourier 变换为 $(2\pi)^{m/2}\delta$. 由 \eqref{eq:fourier-frequency-derivatives}, \eqref{eq:fourier-spatial-derivatives}, 现在可知多项式的 Fourier 变换是 delta 函数导数的线性组合, 反之亦然.
\end{Example}

\begin{Example}
\label{ex:fourier-exponential-periodic}
与 Example~\ref{ex:fourier-constant} 中类似的计算表明, $\exp(i\bm{\eta}\cdot\bm{x})$ (视为 $\bm{x}$ 的函数) 的 Fourier 变换为 $(2\pi)^{m/2}\delta(\bm{\xi}-\bm{\eta})$. 若 $f$ 是由 Fourier 级数给出的周期分布
\begin{equation}\label{eq:periodic-distribution-fourier-series}
f(x)=\sum_{n=-\infty}^{\infty}c_ne^{inx},
\end{equation}
则得到
\begin{equation}\label{eq:periodic-distribution-fourier-point-masses}
\mathcal{F}[f](\xi)=\sqrt{2\pi}\sum_{n=-\infty}^{\infty}c_n\delta(\xi-n).
\end{equation}

\end{Example}

\begin{Example}
\label{ex:fourier-compact-distribution-entire}
设 $f$ 为紧支集分布. 那么可对任意 $\phi\in C^\infty(\mathbb{R}^m)$ 定义 $(f,\phi)$: 令 $(f,\phi)=(f,\phi_0)$, 其中 $\phi_0$ 为 $\mathcal{D}(\mathbb{R}^m)$ 中任意在 $f$ 的支集的某邻域内与 $\phi$ 一致的元素. 由支集的定义可知, 这一定义不依赖于 $\phi_0$ 的选择. 特别地, 这将 $f$ 定义为 $\mathcal{S}^{\prime}(\mathbb{R}^m)$ 的一个元素. 现在断言 $\mathcal{F}[f]$ 是函数
\begin{equation}\label{eq:compact-distribution-fourier-entire-formula}
\mathcal{F}[f](\bm{\xi})=(2\pi)^{-m/2}(\overline{f(\bm{x})},e^{-i\bm{\xi}\cdot\bm{x}}).
\end{equation}
这里 $(\overline f,\phi)$ 定义为 $(f,\overline\phi)$ 的复共轭. 为验证这一断言, 必须证明对任意 $\phi\in\mathcal{S}(\mathbb{R}^m)$, 有
\begin{equation}\label{eq:compact-distribution-fourier-integral-interchange}
\begin{aligned}
&(2\pi)^{-m/2}\int_{\mathbb{R}^m}(f(\bm{x}),e^{i\bm{\xi}\cdot\bm{x}})\phi(\bm{\xi})\,\mathrm{d}\bm{\xi}\\
&\quad=(f,\mathcal{F}^{-1}[\phi])\\
&\quad=(2\pi)^{-m/2}\left(f,\int_{\mathbb{R}^m}e^{i\bm{\xi}\cdot\bm{x}}\phi(\bm{\xi})\,\mathrm{d}\bm{\xi}\right).
\end{aligned}
\end{equation}
即必须说明将 $f$ 移入积分号内是合理的, 这可通过以有限和逼近积分, 按通常方式完成. 注意, \eqref{eq:compact-distribution-fourier-entire-formula} 定义了 $\bm{\xi}\in\mathbb{C}^m$ 的整函数. 紧支集分布具有有限阶这一事实 (Lemma~\ref{lem:distribution-local-finite-order}) 意味着对实自变量, 这一函数具有多项式增长.
\end{Example}
